% Columna de material en caída entre la salida del tornillo y el material ya
% pesado, según el modelo de Bukreev et al. (2014). Sin escala.
\begin{tikzpicture}[
    >=Latex,
    font=\small,
    nota/.style={font=\footnotesize, align=center},
    cota/.style={{Latex[width=4pt,length=4pt]}-{Latex[width=4pt,length=4pt]}},
    celda/.style={draw, thick, fill=gray!30}
]

% Tornillo con salida hacia abajo
\draw[thick] (0,6.0) rectangle (3.5,6.6);
\foreach \x in {0.2,0.6,1.0,1.4,1.8,2.2,2.6}{\draw (\x,6.0) -- (\x+0.3,6.6);}
\draw[thick] (2.9,6.0) -- (2.9,5.8) (3.4,6.0) -- (3.4,5.8);
\node[nota, anchor=south] at (1.75,6.65) {Tornillo};

% Columna en caída
\fill[gray!45] (2.9,5.8) rectangle (3.4,2.0);
\draw[densely dashed] (2.9,5.8) -- (2.9,2.0) (3.4,5.8) -- (3.4,2.0);
\draw[gray] (2.9,4.6) -- (0.9,4.6);
\node[nota, anchor=east] at (0.9,4.6) {Columna\\en caída};

% Recipiente y material pesado
\fill[gray!20] (1.5,0.6) rectangle (4.8,2.0);
\draw[thick] (1.5,4.0) -- (1.5,0.6) -- (4.8,0.6) -- (4.8,4.0);
\draw[gray] (1.6,1.3) -- (0.9,1.3);
\node[nota, anchor=east] at (0.9,1.3) {Material\\pesado};

% Celda de carga
\draw[celda] (2.8,0.3) rectangle (3.5,0.6);
\draw[thick] (2.5,0.3) -- (3.8,0.3);
\node[nota, anchor=north] at (3.15,0.25) {Celda de carga};

% Cotas
\draw[densely dotted, gray] (3.4,6.0) -- (6.9,6.0);
\draw[densely dotted, gray] (4.8,4.0) -- (6.9,4.0);
\draw[densely dotted, gray] (4.8,2.0) -- (6.0,2.0);
\draw[densely dotted, gray] (4.8,0.6) -- (6.9,0.6);
\draw[cota] (5.6,6.0) -- (5.6,4.0) node[midway, right] {$h_0$};
\draw[cota] (6.6,4.0) -- (6.6,0.6) node[midway, right] {$h_b$};
\draw[cota] (5.6,2.0) -- (5.6,0.6) node[midway, right] {$h_m$};
\draw[cota] (4.2,5.8) -- (4.2,2.0) node[midway, right] {$h_c$};

\end{tikzpicture}
